Fix \(t\in\{0,1\}\). By \(\mathrm{ass:gaussian\text{-}errors}\), there are a regular conditional law \(K_{t,P}(z,du)\) of \(\varepsilon_t=Y-g_{t,P}(X)\) given \(X=z\) and a \(P_X\)-null set \(N_{t,P}\) such that all three clauses of that assumption hold simultaneously for every \(z\in\mathcal X_{t,P}\setminus N_{t,P}\). Fix such a \(z\), and abbreviate \(K(du)=K_{t,P}(z,du)\). The first-moment and centering clauses give directly
\[
 \int_{\mathbb R}|u|K(du)<\infty,
 \qquad
 \int_{\mathbb R}uK(du)=0.
 \tag{1}
\]
Thus \(\varepsilon_t\) is conditionally integrable and \(\mathbb E_P[\varepsilon_t\mid X=z]=0\).

The assumed moment-generating-function inequality is already the claimed conditional inequality: for every \(\lambda\in\mathbb R\),
\[
 \mathbb E_P[\exp(\lambda\varepsilon_t)\mid X=z]
 =\int_{\mathbb R}e^{\lambda u}K(du)
 \le e^{\sigma^2\lambda^2/2}.
 \tag{2}
\]
It remains to derive the displayed numerical first-moment bound. Fix \(s>0\). For every \(\lambda>0\), the pointwise inequality
\[
 \mathbf 1\{u\ge s\}\le e^{\lambda(u-s)}
\]
and (2) imply
\[
 K([s,\infty))
 \le \exp\!\left(-\lambda s+\frac{\sigma^2\lambda^2}{2}\right).
 \tag{3}
\]
Because \(\sigma>0\), the choice \(\lambda=s/\sigma^2\) is valid and minimizes the quadratic exponent, so
\[
 K([s,\infty))\le e^{-s^2/(2\sigma^2)}.
 \tag{4}
\]
Applying (2) with \(-\lambda\), or equivalently applying the same argument to \(-u\), gives
\[
 K(({-\infty},-s])\le e^{-s^2/(2\sigma^2)}.
 \tag{5}
\]
For each \(u\in\mathbb R\), \(|u|=\int_0^\infty\mathbf 1\{s<|u|\}\,ds\). Tonelli's identity for this nonnegative integrand, followed by (4)--(5), therefore yields
\[
 \begin{aligned}
 \mathbb E_P[|\varepsilon_t|\mid X=z]
 &=\int_{\mathbb R}|u|K(du)\\
 &=\int_0^\infty K\{|u|>s\}\,ds\\
 &\le 2\int_0^\infty e^{-s^2/(2\sigma^2)}\,ds
 =\sigma\sqrt{2\pi}.
 \end{aligned}
 \tag{6}
\]
Finally, \(g_{t,P}(z)\in\mathbb R\) by its declared signature. Hence (1) makes \(Y=g_{t,P}(z)+\varepsilon_t\) integrable conditionally at this \(z\), and conditional centering gives
\[
 \mathbb E_P[Y\mid X=z]
 =g_{t,P}(z)+\mathbb E_P[\varepsilon_t\mid X=z]
 =g_{t,P}(z).
 \tag{7}
\]
The exceptional set \(N_{t,P}\) is \(P_X\)-null, and \(t\) was arbitrary. Equations (1), (2), (6), and (7) prove every assertion for each \(t\in\{0,1\}\) and for \(P_X\)-almost every \(z\in\mathcal X_{t,P}\).